% ---------------------------------------------------------
\section{Limits Involving Infinity}

% ---------------------------------------------------------
\section*{Infinite Limits}

\begin{propositionbox}{Definition (Limits of Infinity)}
Let $I$ be an open interval containing $c$, and let $f$ be defined
on $I$, except possibly at $c$.
\begin{itemize}
  \item The limit of $f(x)$ as $x$ approaches $c$ is \emph{infinity},
        denoted
        \[
          \lim_{x \to c} f(x) = \infty,
        \]
        if given any $N > 0$, there exists $\delta > 0$ such that
        for all $x \in I$ with $x \neq c$, if $|x - c| < \delta$
        then $f(x) > N$.

  \item The limit of $f(x)$ as $x$ approaches $c$ is
        \emph{negative infinity}, denoted
        \[
          \lim_{x \to c} f(x) = -\infty,
        \]
        if given any $N < 0$, there exists $\delta > 0$ such that
        for all $x \in I$ with $x \neq c$, if $|x - c| < \delta$
        then $f(x) < N$.
\end{itemize}
\end{propositionbox}

\begin{remark*}[Intuition]
$N$ plays the role that $\varepsilon$ plays in the standard
definition: the adversary names an arbitrarily large threshold,
and the response is a $\delta$-neighbourhood of $c$ on which the
function exceeds it.
Infinity is not a value being approached --- it is a description
of unbounded growth.
\end{remark*}

% ---------------------------------------------------------
\section*{One-Sided Limits of Infinity}

\begin{propositionbox}{Definition (One-Sided Limits of Infinity)}
\textbf{Left-hand limit.}
Let $f$ be defined on $(a, c)$ for some $a < c$.
The left-hand limit of $f$ at $c$ is infinity,
\[
  \lim_{x \to c^-} f(x) = \infty,
\]
if given any $N > 0$, there exists $\delta > 0$ such that for all
$a < x < c$, if $|x - c| < \delta$ then $f(x) > N$.

\medskip
\textbf{Right-hand limit.}
Let $f$ be defined on $(c, b)$ for some $b > c$.
The right-hand limit of $f$ at $c$ is infinity,
\[
  \lim_{x \to c^+} f(x) = \infty,
\]
if given any $N > 0$, there exists $\delta > 0$ such that for all
$c < x < b$, if $|x - c| < \delta$ then $f(x) > N$.

\medskip
The left- and right-hand limits of $-\infty$ are defined analogously,
replacing $f(x) > N$ with $f(x) < N$ for $N < 0$.
\end{propositionbox}

% ---------------------------------------------------------
\section*{Vertical Asymptotes}

\begin{propositionbox}{Definition (Vertical Asymptote)}
Let $I$ be an interval that either contains $c$ or has $c$ as an
endpoint, and let $f$ be defined on $I$, except possibly at $c$.

If the limit of $f(x)$ as $x$ approaches $c$ from either the left
or the right (or both) is $\infty$ or $-\infty$, then the line
$x = c$ is a \emph{vertical asymptote} of $f$.
\end{propositionbox}

\begin{remark*}
The standard example is $f(x) = 1/x$, which has a vertical asymptote
at $x = 0$: $\lim_{x \to 0^-} 1/x = -\infty$ and
$\lim_{x \to 0^+} 1/x = +\infty$.
The two one-sided limits need not agree --- or even have the same
sign --- for $x = c$ to qualify as a vertical asymptote.
\end{remark*}

% ---------------------------------------------------------
\section*{Indeterminate Forms}

\begin{propositionbox}{Definition (Indeterminate Forms)}
An \emph{indeterminate form} arises when substituting the limit
value into an expression produces a form whose limiting value
cannot be determined from the form alone.
The standard indeterminate forms are:
\[
  \frac{0}{0}, \quad
  \frac{\infty}{\infty}, \quad
  0 \cdot \infty, \quad
  \infty - \infty, \quad
  0^0, \quad
  1^\infty, \quad
  \infty^0.
\]
\end{propositionbox}

\begin{remark*}
An indeterminate form is not an error --- it is a signal that
more work is required.
The form $0/0$ is the most common: it appears whenever a common
factor vanishes at the limit point (handled by cancellation)
or in derivatives via the difference quotient.
The forms $0/0$ and $\infty/\infty$ are handled by
L'H\^opital's Rule; the others are typically reduced to one
of these two by algebraic manipulation.
\end{remark*}

\begin{remark*}[Common error]
The forms $0/0$ and $\infty/\infty$ are indeterminate, but $c/0$
for $c \neq 0$ is \emph{not} indeterminate --- it indicates an
infinite limit (a vertical asymptote), not an unresolved
expression.
Similarly, $\infty + \infty = \infty$ is determinate; only
$\infty - \infty$ is not.
\end{remark*}

% ---------------------------------------------------------
\section*{Limits at Infinity and Horizontal Asymptotes}

\begin{propositionbox}{Definition (Limits at Infinity and Horizontal Asymptotes)}
Let $L$ be a real number.
\begin{enumerate}
  \item Let $f$ be defined on $(a, \infty)$ for some $a$.
        The \emph{limit of $f$ at infinity} is $L$, denoted
        $\lim_{x \to \infty} f(x) = L$, if for every
        $\varepsilon > 0$ there exists $M > a$ such that if
        $x > M$ then $|f(x) - L| < \varepsilon$.

  \item Let $f$ be defined on $(-\infty, b)$ for some $b$.
        The \emph{limit of $f$ at negative infinity} is $L$,
        denoted $\lim_{x \to -\infty} f(x) = L$, if for every
        $\varepsilon > 0$ there exists $M < b$ such that if
        $x < M$ then $|f(x) - L| < \varepsilon$.

  \item If $\lim_{x \to \infty} f(x) = L$ or
        $\lim_{x \to -\infty} f(x) = L$, the line $y = L$ is a
        \emph{horizontal asymptote} of $f$.
\end{enumerate}
\end{propositionbox}

\begin{remark*}
A function can have at most two horizontal asymptotes (one as
$x \to +\infty$ and one as $x \to -\infty$), and they need not
be equal.
For example, $f(x) = \arctan(x)$ has horizontal asymptotes at
$y = \pi/2$ and $y = -\pi/2$.
\end{remark*}

% ---------------------------------------------------------
\section*{Limits of Rational Functions at Infinity}

\begin{theorembox}{Theorem (Limits of Rational Functions at Infinity)}
Let $f(x)$ be a rational function of the form
\[
  f(x) = \frac{a_n x^n + a_{n-1} x^{n-1} + \cdots + a_1 x + a_0}
              {b_m x^m + b_{m-1} x^{m-1} + \cdots + b_1 x + b_0},
\]
where $m, n$ are positive integers and $a_n, b_m \neq 0$. Then:
\begin{enumerate}
  \item If $n = m$, then
        $\displaystyle\lim_{x \to \infty} f(x)
         = \lim_{x \to -\infty} f(x) = \dfrac{a_n}{b_m}$.

  \item If $n < m$, then
        $\displaystyle\lim_{x \to \infty} f(x)
         = \lim_{x \to -\infty} f(x) = 0$.

  \item If $n > m$, then $\lim_{x \to \infty} f(x)$ and
        $\lim_{x \to -\infty} f(x)$ are both infinite.
\end{enumerate}
\end{theorembox}

\begin{remark*}
The rule follows from dividing numerator and denominator by $x^m$
and observing that all terms with negative powers of $x$ vanish
as $x \to \pm\infty$, leaving only the ratio $a_n x^{n-m} / b_m$.
When $n = m$ this ratio is the constant $a_n/b_m$; when $n < m$
it tends to zero; when $n > m$ it grows without bound.
\end{remark*}